
\documentclass{article}
\usepackage{colortbl}
\usepackage{makecell}
\usepackage{multirow}
\usepackage{supertabular}

\begin{document}

\newcounter{utterance}

\centering \large Interaction Transcript for game `dond', experiment `semi\_en', episode 13 with Qwen3.5{-}4B.
\vspace{24pt}

{ \footnotesize  \setcounter{utterance}{1}
\setlength{\tabcolsep}{0pt}
\begin{supertabular}{c@{$\;$}|p{.15\linewidth}@{}p{.15\linewidth}p{.15\linewidth}p{.15\linewidth}p{.15\linewidth}p{.15\linewidth}}
    \# & $\;$A & \multicolumn{4}{c}{Game Master} & $\;\:$B\\
    \hline

    \theutterance \stepcounter{utterance}  
    & & \multicolumn{4}{p{0.6\linewidth}}{
        \cellcolor[rgb]{0.9,0.9,0.9}{
            \makecell[{{p{\linewidth}}}]{
                \texttt{\tiny{[P1$\langle$GM]}}
                \texttt{You are playing a negotiation game in which you have to agree on how to divide a set of items among you and another player.} \\
\\ 
\texttt{Rules:} \\
\texttt{(a) You and the other player are given a set of items. Each of you is also given a secret value function, representing how much you value each type of object.} \\
\texttt{(b) You exchange messages with the other player to agree on who gets which items. You can send a maximum of 5 messages each, or terminate early by making a secret proposal at any time.} \\
\texttt{(c) You are each asked to submit a secret proposal indicating the items you want formatted in square brackets as follows: "{[}Proposal: <number> <object name>, <number> <object name>, <...>{]}"} \\
\texttt{(d) If your proposals are complementary, i.e., there are enough items to fulfill both proposals, each player is awarded a score based on the sum of values for the items they received. Otherwise, both of you get zero points.} \\
\texttt{(e) Your goal is to maximize the score you receive.} \\
\\ 
\texttt{Let us start.} \\
\\ 
\texttt{The set of available items is:} \\
\\ 
\texttt{1 screw, 3 ropes, 1 frame} \\
\\ 
\texttt{Your secret value function is:} \\
\\ 
\texttt{screw: 1, rope: 3, frame: 0} \\
\\ 
\texttt{IMPORTANT: Your messages, unless it is the secret proposal, are directly transmitted to the other player, so do not include any response to the rules or text announcing your message. To make a secret proposal, use the indicated format. Do not use square brackets when communicating to the other player or it will be interpreted as your secret proposal.} \\
            }
        }
    }
    & \\ \\

    \theutterance \stepcounter{utterance}  
    & \multicolumn{4}{p{0.6\linewidth}}{
        \cellcolor[rgb]{0.9,0.9,0.9}{
            \makecell[{{p{\linewidth}}}]{
                \texttt{\tiny{[P1$\rangle$GM]}}
                \texttt{I will propose a division that maximizes my score while ensuring the other player also receives a positive value, making the proposals complementary.} \\
\\ 
\texttt{My valuation:} \\
\texttt{{-} 1 screw: 1} \\
\texttt{{-} 3 ropes: 3 (total)} \\
\texttt{{-} 1 frame: 0} \\
\\ 
\texttt{Total value = 4.} \\
\\ 
\texttt{Strategy:} \\
\texttt{{-} I value the ropes most (3 points). I should get all 3 ropes.} \\
\texttt{{-} The other player gets the screw (1 point) and the frame (0 points). Total = 1.} \\
\texttt{{-} This is complementary (items are fully distributed).} \\
\texttt{{-} My score: 3.} \\
\texttt{{-} Alternative: If I take the screw and frame (value 1) and give ropes to the other player (value 3), my score is 1.} \\
\texttt{{-} Alternative: If I take the screw and one rope (value 4) and give the other two ropes and frame to the other player (value 6), my score is 4. Wait, let's re{-}evaluate.} \\
\texttt{  {-} Option A: I get 3 ropes (3 pts). Other gets screw + frame (1 pt). Total 4. My score 3.} \\
\texttt{  {-} Option B: I get screw + 1 rope (1+3=4 pts). Other gets 2 ropes + frame (6+0=6 pts). Total 10? No, total value is fixed.} \\
\texttt{    {-} Items: 1 screw, 3 ropes, 1 frame.} \\
\texttt{    {-} My val: S=1, R=3, F=0. Total = 1 + 3*3 + 0 = 10? No, "3 ropes" means there are 3 distinct rope items.} \\
\texttt{    {-} Let's re{-}read carefully: "1 screw, 3 ropes, 1 frame".} \\
\texttt{    {-} My value function: screw: 1, rope: 3, frame: 0.} \\
\texttt{    {-} Total value for me = 1 (screw) + 3*3 (three ropes) + 0 (frame) = 10.} \\
\texttt{    {-} If I take all 3 ropes, I get 3*3 = 9.} \\
\texttt{    {-} Remaining items: 1 screw, 1 frame.} \\
\texttt{    {-} Other player gets screw (1) + frame (0) = 1.} \\
\texttt{    {-} Sum = 10.} \\
\texttt{    {-} If I take screw + 1 rope: Value = 1 + 3 = 4.} \\
\texttt{    {-} Remaining: 2 ropes, 1 frame.} \\
\texttt{    {-} Other player gets 2 ropes + frame: Value = 2*3 + 0 = 6.} \\
\texttt{    {-} Sum = 10.} \\
\texttt{    {-} If I take screw + 2 ropes: Value = 1 + 6 = 7.} \\
\texttt{    {-} Remaining: 1 rope} \\
            }
        }
    }
    & & \\ \\

    \theutterance \stepcounter{utterance}  
    & & & \multicolumn{2}{p{0.3\linewidth}}{
        \cellcolor[rgb]{0.9,0.9,0.9}{
            \makecell[{{p{\linewidth}}}]{
                \texttt{\tiny{[GM$|$GM]}}
                \texttt{continue} \\
            }
        }
    }
    & & \\ \\

    \theutterance \stepcounter{utterance}  
    & & & \multicolumn{2}{p{0.3\linewidth}}{
        \cellcolor[rgb]{0.9,0.9,0.9}{
            \makecell[{{p{\linewidth}}}]{
                \texttt{\tiny{[GM$|$GM]}}
                \texttt{I will propose a division that maximizes my score while ensuring the other player also receives a positive value, making the proposals complementary.} \\
\\ 
\texttt{My valuation:} \\
\texttt{{-} 1 screw: 1} \\
\texttt{{-} 3 ropes: 3 (total)} \\
\texttt{{-} 1 frame: 0} \\
\\ 
\texttt{Total value = 4.} \\
\\ 
\texttt{Strategy:} \\
\texttt{{-} I value the ropes most (3 points). I should get all 3 ropes.} \\
\texttt{{-} The other player gets the screw (1 point) and the frame (0 points). Total = 1.} \\
\texttt{{-} This is complementary (items are fully distributed).} \\
\texttt{{-} My score: 3.} \\
\texttt{{-} Alternative: If I take the screw and frame (value 1) and give ropes to the other player (value 3), my score is 1.} \\
\texttt{{-} Alternative: If I take the screw and one rope (value 4) and give the other two ropes and frame to the other player (value 6), my score is 4. Wait, let's re{-}evaluate.} \\
\texttt{  {-} Option A: I get 3 ropes (3 pts). Other gets screw + frame (1 pt). Total 4. My score 3.} \\
\texttt{  {-} Option B: I get screw + 1 rope (1+3=4 pts). Other gets 2 ropes + frame (6+0=6 pts). Total 10? No, total value is fixed.} \\
\texttt{    {-} Items: 1 screw, 3 ropes, 1 frame.} \\
\texttt{    {-} My val: S=1, R=3, F=0. Total = 1 + 3*3 + 0 = 10? No, "3 ropes" means there are 3 distinct rope items.} \\
\texttt{    {-} Let's re{-}read carefully: "1 screw, 3 ropes, 1 frame".} \\
\texttt{    {-} My value function: screw: 1, rope: 3, frame: 0.} \\
\texttt{    {-} Total value for me = 1 (screw) + 3*3 (three ropes) + 0 (frame) = 10.} \\
\texttt{    {-} If I take all 3 ropes, I get 3*3 = 9.} \\
\texttt{    {-} Remaining items: 1 screw, 1 frame.} \\
\texttt{    {-} Other player gets screw (1) + frame (0) = 1.} \\
\texttt{    {-} Sum = 10.} \\
\texttt{    {-} If I take screw + 1 rope: Value = 1 + 3 = 4.} \\
\texttt{    {-} Remaining: 2 ropes, 1 frame.} \\
\texttt{    {-} Other player gets 2 ropes + frame: Value = 2*3 + 0 = 6.} \\
\texttt{    {-} Sum = 10.} \\
\texttt{    {-} If I take screw + 2 ropes: Value = 1 + 6 = 7.} \\
\texttt{    {-} Remaining: 1 rope} \\
            }
        }
    }
    & & \\ \\

    \theutterance \stepcounter{utterance}  
    & & \multicolumn{4}{p{0.6\linewidth}}{
        \cellcolor[rgb]{0.9,0.9,0.9}{
            \makecell[{{p{\linewidth}}}]{
                \texttt{\tiny{[GM$\rangle$P2]}}
                \texttt{You are playing a negotiation game in which you have to agree on how to divide a set of items among you and another player.} \\
\\ 
\texttt{Rules:} \\
\texttt{(a) You and the other player are given a set of items. Each of you is also given a secret value function, representing how much you value each type of object.} \\
\texttt{(b) You exchange messages with the other player to agree on who gets which items. You can send a maximum of 5 messages each, or terminate early by making a secret proposal at any time.} \\
\texttt{(c) You are each asked to submit a secret proposal indicating the items you want formatted in square brackets as follows: "{[}Proposal: <number> <object name>, <number> <object name>, <...>{]}"} \\
\texttt{(d) If your proposals are complementary, i.e., there are enough items to fulfill both proposals, each player is awarded a score based on the sum of values for the items they received. Otherwise, both of you get zero points.} \\
\texttt{(e) Your goal is to maximize the score you receive.} \\
\\ 
\texttt{Let us start.} \\
\\ 
\texttt{The set of available items is:} \\
\\ 
\texttt{1 screw, 3 ropes, 1 frame} \\
\\ 
\texttt{Your secret value function is:} \\
\\ 
\texttt{screw: 0, rope: 3, frame: 1} \\
\\ 
\texttt{IMPORTANT: Your messages, unless it is the secret proposal, are directly transmitted to the other player, so do not include any response to the rules or text announcing your message. To make a secret proposal, use the indicated format. Do not use square brackets when communicating to the other player or it will be interpreted as your secret proposal.} \\
\\ 
\\ 
\texttt{I will propose a division that maximizes my score while ensuring the other player also receives a positive value, making the proposals complementary.} \\
\\ 
\texttt{My valuation:} \\
\texttt{{-} 1 screw: 1} \\
\texttt{{-} 3 ropes: 3 (total)} \\
\texttt{{-} 1 frame: 0} \\
\\ 
\texttt{Total value = 4.} \\
\\ 
\texttt{Strategy:} \\
\texttt{{-} I value the ropes most (3 points). I should get all 3 ropes.} \\
\texttt{{-} The other player gets the screw (1 point) and the frame (0 points). Total = 1.} \\
\texttt{{-} This is complementary (items are fully distributed).} \\
\texttt{{-} My score: 3.} \\
\texttt{{-} Alternative: If I take the screw and frame (value 1) and give ropes to the other player (value 3), my score is 1.} \\
\texttt{{-} Alternative: If I take the screw and one rope (value 4) and give the other two ropes and frame to the other player (value 6), my score is 4. Wait, let's re{-}evaluate.} \\
\texttt{  {-} Option A: I get 3 ropes (3 pts). Other gets screw + frame (1 pt). Total 4. My score 3.} \\
\texttt{  {-} Option B: I get screw + 1 rope (1+3=4 pts). Other gets 2 ropes + frame (6+0=6 pts). Total 10? No, total value is fixed.} \\
\texttt{    {-} Items: 1 screw, 3 ropes, 1 frame.} \\
\texttt{    {-} My val: S=1, R=3, F=0. Total = 1 + 3*3 + 0 = 10? No, "3 ropes" means there are 3 distinct rope items.} \\
\texttt{    {-} Let's re{-}read carefully: "1 screw, 3 ropes, 1 frame".} \\
\texttt{    {-} My value function: screw: 1, rope: 3, frame: 0.} \\
\texttt{    {-} Total value for me = 1 (screw) + 3*3 (three ropes) + 0 (frame) = 10.} \\
\texttt{    {-} If I take all 3 ropes, I get 3*3 = 9.} \\
\texttt{    {-} Remaining items: 1 screw, 1 frame.} \\
\texttt{    {-} Other player gets screw (1) + frame (0) = 1.} \\
\texttt{    {-} Sum = 10.} \\
\texttt{    {-} If I take screw + 1 rope: Value = 1 + 3 = 4.} \\
\texttt{    {-} Remaining: 2 ropes, 1 frame.} \\
\texttt{    {-} Other player gets 2 ropes + frame: Value = 2*3 + 0 = 6.} \\
\texttt{    {-} Sum = 10.} \\
\texttt{    {-} If I take screw + 2 ropes: Value = 1 + 6 = 7.} \\
\texttt{    {-} Remaining: 1 rope} \\
            }
        }
    }
    & \\ \\

    \theutterance \stepcounter{utterance}  
    & & & \multicolumn{4}{p{0.6\linewidth}}{
        \cellcolor[rgb]{0.9,0.9,0.9}{
            \makecell[{{p{\linewidth}}}]{
                \texttt{\tiny{[GM$\langle$P2]}}
                \texttt{{[}Proposal: 3 ropes{]}} \\
            }
        }
    }
     \\ \\

    \theutterance \stepcounter{utterance}  
    & & & \multicolumn{2}{p{0.3\linewidth}}{
        \cellcolor[rgb]{0.9,0.9,0.9}{
            \makecell[{{p{\linewidth}}}]{
                \texttt{\tiny{[GM$|$GM]}}
                \texttt{proposal} \\
            }
        }
    }
    & & \\ \\

    \theutterance \stepcounter{utterance}  
    & & & \multicolumn{2}{p{0.3\linewidth}}{
        \cellcolor[rgb]{0.9,0.9,0.9}{
            \makecell[{{p{\linewidth}}}]{
                \texttt{\tiny{[GM$|$GM]}}
                [0, 3, 0]
            }
        }
    }
    & & \\ \\

    \theutterance \stepcounter{utterance}  
    & & \multicolumn{4}{p{0.6\linewidth}}{
        \cellcolor[rgb]{0.9,0.9,0.9}{
            \makecell[{{p{\linewidth}}}]{
                \texttt{\tiny{[P1$\langle$GM]}}
                \texttt{The other player has submitted a secret proposal. It is now your turn to submit a proposal of your own.} \\
            }
        }
    }
    & \\ \\

    \theutterance \stepcounter{utterance}  
    & \multicolumn{4}{p{0.6\linewidth}}{
        \cellcolor[rgb]{0.9,0.9,0.9}{
            \makecell[{{p{\linewidth}}}]{
                \texttt{\tiny{[P1$\rangle$GM]}}
                \texttt{{[}Proposal: 3 ropes{]}} \\
            }
        }
    }
    & & \\ \\

    \theutterance \stepcounter{utterance}  
    & & & \multicolumn{2}{p{0.3\linewidth}}{
        \cellcolor[rgb]{0.9,0.9,0.9}{
            \makecell[{{p{\linewidth}}}]{
                \texttt{\tiny{[GM$|$GM]}}
                \texttt{proposal} \\
            }
        }
    }
    & & \\ \\

    \theutterance \stepcounter{utterance}  
    & & & \multicolumn{2}{p{0.3\linewidth}}{
        \cellcolor[rgb]{0.9,0.9,0.9}{
            \makecell[{{p{\linewidth}}}]{
                \texttt{\tiny{[GM$|$GM]}}
                [0, 3, 0]
            }
        }
    }
    & & \\ \\

    \theutterance \stepcounter{utterance}  
    & & & \multicolumn{2}{p{0.3\linewidth}}{
        \cellcolor[rgb]{0.9,0.9,0.9}{
            \makecell[{{p{\linewidth}}}]{
                \texttt{\tiny{[GM$|$GM]}}
                [[0, 3, 0], [0, 3, 0]]
            }
        }
    }
    & & \\ \\

\end{supertabular}
}

\end{document}
